\section{Related Work}

\label{sec:related}

\subsection{General-Purpose PaaS}

Heroku~\cite{heroku2007} pioneered the PaaS model with git-push deployment. Vercel~\cite{vercel2020} specializes in frontend and serverless deployments. Railway~\cite{railway2021} provides database-friendly PaaS. Fly.io~\cite{flyio2020} offers edge computing with container support. Render~\cite{render2019} provides managed infrastructure. None of these treat AI workloads as first-class citizens.

\subsection{AI-Specific Infrastructure}

Modal~\cite{modal2023} provides serverless GPU compute for AI workloads. Replicate~\cite{replicate2022} offers model hosting with API generation. Baseten~\cite{baseten2022} specializes in ML model deployment. Together AI~\cite{together2023} provides inference API with custom model support. Hanzo Platform differs by providing a full application platform (not just model serving) with integrated databases, secrets, and frontend hosting.

\subsection{Kubernetes-Based PaaS}

Dokku~\cite{dokku2013} provides Heroku-like deployment on single servers. KubeSphere~\cite{kubesphere2020} offers Kubernetes-based application management. Dokploy~\cite{dokploy2024} provides self-hosted PaaS on Kubernetes. Hanzo Platform builds on Dokploy's foundation with AI-specific extensions for GPU management, model caching, and inference autoscaling.

\subsection{ML Deployment Platforms}

MLflow~\cite{mlflow2018} provides ML lifecycle management. Seldon Core~\cite{seldon2019} offers ML model serving on Kubernetes. KServe~\cite{kserve2021} provides serverless inference. BentoML~\cite{bentoml2022} packages ML models for deployment. These tools focus on model serving; Hanzo Platform encompasses the full application stack.

